\section{Introduction}
A vast library of papyrus scrolls was carbonized at Herculaneum during the 79 AD eruption of Mount Vesuvius. Synchrotron micro-CT preserves their internal geometry, but the ink becomes readable only after the correct papyrus surface has been found in three dimensions and {\em virtually unwrapped} to a plane. The useful computational object is therefore not a closed volume boundary. It is one consistently oriented side of an open sheet, with adjacent turns kept separate, genuine holes distinguished from meshing holes, and enough metric fidelity that fibres and writing remain continuous after flattening.

Automatic segmenters based on U-Net architectures~\cite{ronneberger2015u} provide strong voxel-level evidence, but not that object. Their predictions can be several voxels thick, contain gaps and solid blocks, or join adjacent wraps through narrow and broad {\em fusion bridges}. Standard reconstruction tools compound the mismatch: marching cubes~\cite{lorensen1998marching} and Poisson reconstruction~\cite{kazhdan2006poisson} seek a two-sided, watertight boundary. Meshing a scroll requires an algorithm that outputs a single-sided, open, and approximately developable manifold. Moreover, results that look plausible when rendered may still contain small defects; a single short edge connecting two surfaces a full turn apart can invalidate every downstream parameterization. This in turn introduces a second challenge: building a robust parameterization of the meshed scroll that corrects both for scroll damage, such as sheet delamination and physical damage to the scroll itself, as well as any errors introduced during meshing. Finally, any algorithm for scroll meshing must be parallelizable and be able to operate on small domains within the larger volumetric data, which consists of tens of thousands of cubes of voxel geometry for a single scroll.

ScrollFiesta addresses this challenge, treating meshing, topology, parameterization, and validation as one pipeline while operating on individual cubes of CT data and predictions. Given raw CT data and a matching prediction (generated by the existing nnU-Net pipelines), it constructs single-sided meshes of open sheets using a pipeline based on a compact-support moving least squares fit (MLS) combined with a guarded ball-pivoting meshing. It then separates bridged or overlapping layers, caused either by scroll damage or errors in the original prediction; fills clean boundary loops; and severs short topological handles. The resulting mesh sheets are remeshed with a domain-specific mesh level of detail algorithm, based around restricted Voronoi diagrams (RVD). Per-cube results are then joined into larger surfaces by inter-cube seam welding.

The output is then audited and parameterized according to its topology. ScrollFiesta uses a native slice-domain solve, attempting to fit a global arc-length parameterized scroll to the reconstructed mesh sheets. Simply using winding order produces invalid results, and breaks down in the presence of close together sheets; at the same time, it is also necessary to be able to parameterize the scroll even in the presence of such sheets, scroll damage, or the meshing process itself. We therefore parameterize the scroll into one long strip in UV space using a combined discrete-continuous optimization, in which scroll pieces are optimized using local winding information and then globally fit into the correct striation using a tabu search. We then use this map to identify problems such as delamination, off-sheet prediction, and seam mismatches, and repair them using a local quilting energy.

Verified local patches can be imported into VC3D's native loaders and connectors. Larger surfaces can now be exported as new, provenance-bearing {\em unverified Spiral hints}. We provide command-line tools, a public C ABI and CMake package, and a Python workflow that can stream an OME-Zarr region directly from S3.

Our technical contributions are as follows. First, we provide a local, single-sided reconstruction and topology-repair pipeline designed specifically for segmented scroll volumes. Second, we provide a CVT/RVD based level-of-detail and seam-welding scheme for scrolls with fail-closed topology gates. Third, we provide a complementary toolkit for patch and whole-scroll parameterization, based around an alternating discrete/continuous atlas unwrapping optimization in which a tabu search owns integer winding and registration owns metric placement. Fourth, we provide algorithms for principled scroll repair after unwrapping, based around quilting energy. Finally, all contributions integrate into the existing VC3D/Villa pipelines used by the Herculaneum Papyrus team.

This paper is a progress report: it describes the complete system currently implemented, including stages that work but are not yet final, and reports failures where the audits show that core topology remains unresolved.

